% !TeX program = pdflatex
% Uprawnienia indywidualne inwestora (art. 300^28 KSH) --- wzór postanowień umowy spółki i dokumentów wykonawczych --- Basilisk Systems P.S.A.
% Odesłania pobierane z psa.aux i basilisk_sha.aux (xr-hyper).
\documentclass[11pt,a4paper]{report}
\usepackage[T1]{fontenc}
\usepackage[utf8]{inputenc}
\usepackage[english,polish]{babel}
\usepackage{xr-hyper}
\usepackage{basilisk_i18n}
\usepackage{lmodern}
\usepackage{microtype}
\usepackage{setspace}
\usepackage{parskip}
\usepackage{needspace}
\usepackage{tabularx}
\externaldocument{psa}[psa.pdf]
\externaldocument{basilisk_sha}[basilisk_sha.pdf]

\BusinessPlanCompany{Basilisk Systems}
\BusinessPlanWatermarkText{Basilisk Systems}
\BusinessPlanWatermarkColor{black!6}
\BusinessPlanWatermarkFontSize{38}{46}
\BusinessPlanWatermarkAngle{38}
\BusinessPlanAuthorsText{Basilisk Systems --- wzory dokumentów}
\BusinessPlanConfidentialText{Projekt --- poufne}
\BusinessPlanFooterText{Basilisk Systems --- uprawnienia indywidualne inwestora (wzór)}
\BusinessPlanVersionText{Wersja 1.0 --- projekt}

\setstretch{1.1}
\setlist[enumerate]{topsep=0.2em,itemsep=0.15em,parsep=0.05em,leftmargin=2.0em}
\setlength{\parindent}{0pt}
\setlength{\parskip}{0.5em}
\setlength{\emergencystretch}{4em}
\renewcommand{\arraystretch}{1.2}

\newcounter{paragraf}
\newcommand{\paragraf}[1]{%
  \refstepcounter{paragraf}%
  \Needspace{5\baselineskip}%
  \subsection*{#1}%
  \addcontentsline{toc}{section}{#1}%
}
\newenvironment{ContractClause}[1]{\paragraf{#1}}{\par}
\newcommand{\field}[1]{\textbf{[#1]}}
\newcommand{\CompanyFull}{Basilisk Systems prosta spółka akcyjna}
\newcommand{\US}[1]{§~#1 umowy spółki}
\newcommand{\UA}[1]{§~#1 umowy akcjonariuszy}

\hypersetup{
  pdftitle={Basilisk Systems P.S.A. --- uprawnienia indywidualne inwestora (art. 300^28 KSH)},
  pdfauthor={Basilisk Systems}
}

\begin{document}

\BusinessPlanTitlePage
  {Basilisk Systems}
  {Wersja 1.0 --- projekt}
  {Uprawnienia indywidualne inwestora}
  {Wzór postanowień umowy spółki na podstawie art.~300$^{28}$ KSH}
  {Dyrektor albo obserwator, prawa informacyjne i zgody inwestora}
  {1.0}

\tableofcontents
\newpage

\chapter*{Wprowadzenie}
\addcontentsline{toc}{chapter}{Wprowadzenie}

Umowa spółki może przyznać oznaczonemu akcjonariuszowi uprawnienia indywidualne, w szczególności uprawnienie do powołania lub odwołania członków organów (art.~300$^{28}$ §~1 KSH). Uprawnienia te wygasają najpóźniej z dniem, w którym uprawniony przestaje być akcjonariuszem, chyba że umowa spółki stanowi inaczej (art.~300$^{28}$ §~2 KSH). Wzór realizuje uprawnienia, które warunki Kwalifikowanej Rundy mogą przewidywać zgodnie z \UA{\ref{sha:runda} ust.~2 lit.~c--d}.

\textbf{Zasady stosowania:}
\begin{enumerate}
  \item \textbf{Forma.} Uprawnienia indywidualne muszą wynikać z umowy spółki. Ich wprowadzenie wymaga zmiany umowy uchwałą Walnego Zgromadzenia podjętą większością 75\% wszystkich głosów (\US{\ref{art:wz-zastrzezone} ust.~1}), zaprotokołowaną notarialnie, oraz wpisu do rejestru. Umowa akcjonariuszy może jedynie zobowiązać Strony do głosowania za taką zmianą.
  \item \textbf{Oznaczony akcjonariusz.} Uprawnionego wskazuje się z firmy i numeru rejestrowego. Przyznanie uprawnienia każdorazowemu posiadaczowi akcji określonej serii jest uprzywilejowaniem akcji, a nie uprawnieniem indywidualnym.
  \item \textbf{Ochrona uprawnionego.} Zmiana umowy uszczuplająca uprawnienia indywidualne wymaga zgody uprawnionego (art.~300$^{98}$ §~3 KSH).
  \item \textbf{Wybór wariantów.} Wzór obejmuje cztery warianty: A (dyrektor), B (obserwator), C (prawa informacyjne), D (zgody inwestora). Zwykle stosuje się A albo B, razem z C; wariant D stosuje się tylko w zakresie wynegocjowanym w rundzie.
  \item \textbf{Regulacje.} Prawo powołania dyrektora może wymagać oceny w świetle ustawy o kontroli niektórych inwestycji i zgody organu koncesyjnego (\US{\ref{art:zbywanie} ust.~2}, \UA{\ref{sha:regulacje}}) \field{do weryfikacji}. Dostęp dyrektora inwestora i obserwatora do informacji niejawnych wymaga poświadczenia bezpieczeństwa, a do Informacji Kontrolowanych --- zgody Rady Dyrektorów według polityki bezpieczeństwa informacji.
\end{enumerate}

% ================================================================
\chapter*{Część I --- Postanowienia umowy spółki}
\addcontentsline{toc}{chapter}{Część I --- Postanowienia umowy spółki}

\textbf{Nowy paragraf umowy spółki „Uprawnienia indywidualne”} (po \US{\ref{art:wz-zastrzezone}}):

\begin{ContractClause}{Ust. 1 --- Uprawniony i wygaśnięcie}
\label{upr:uprawniony}
\begin{enumerate}
  \item Uprawnienia indywidualne określone w niniejszym paragrafie przysługują \field{firma inwestora} z siedzibą w \field{miejscowość}, \field{rejestr i numer} (\textbf{„Inwestor”}), na podstawie art.~300$^{28}$ KSH.
  \item Uprawnienia wygasają z dniem, w którym:
  \begin{enumerate}[label=\alph*)]
    \item Inwestor przestaje być akcjonariuszem Spółki;
    \item liczba akcji Inwestora spada poniżej \field{5}\% ogólnej liczby akcji Spółki, z wyłączeniem spadku wynikającego wyłącznie z emisji akcji serii~P i~F; albo
    \item Spółka otrzyma oświadczenie Inwestora o zrzeczeniu się uprawnień.
  \end{enumerate}
  \item \field{Opcjonalnie, do weryfikacji:} Uprawnienia przechodzą na podmiot w całości kontrolowany przez Inwestora albo przez ten sam podmiot zarządzający, który nabędzie wszystkie akcje Inwestora, jeżeli nabywca zawiadomi o tym Spółkę w terminie 14 dni od nabycia.
  \item Rada Dyrektorów zgłasza wygaśnięcie uprawnień do rejestru przedsiębiorców.
\end{enumerate}
\end{ContractClause}

\begin{ContractClause}{Ust. 2 --- Wariant A: dyrektor Inwestora}
\label{upr:dyrektor}
\begin{enumerate}
  \item Inwestor powołuje i odwołuje jednego dyrektora (\textbf{„Dyrektor Inwestora”}) w drodze pisemnego oświadczenia doręczonego Spółce. Oświadczenie jest skuteczne z chwilą doręczenia albo w późniejszym terminie w nim wskazanym.
  \item Dyrektor Inwestora jest dyrektorem niewykonawczym i wchodzi w skład Rady Dyrektorów w granicach liczby dyrektorów określonej w \US{\ref{art:rada} ust.~2}. Pełni funkcję w ramach wspólnej kadencji, a jego mandat wygasa na zasadach z \US{\ref{art:rada} ust.~3}; Inwestor może powołać go ponownie.
  \item Dopóki uprawnienie Inwestora trwa, Walne Zgromadzenie nie powołuje ani nie odwołuje Dyrektora Inwestora. Po wygaśnięciu uprawnienia mandat Dyrektora Inwestora trwa do jego odwołania przez Walne Zgromadzenie.
  \item Jeżeli Inwestor nie powoła Dyrektora Inwestora w terminie 30 dni od wygaśnięcia jego mandatu, miejsce pozostaje nieobsadzone do czasu powołania i nie wpływa na zdolność Rady Dyrektorów do podejmowania uchwał (\US{\ref{art:rada} ust.~5}).
  \item Dyrektor Inwestora może przekazywać Inwestorowi informacje o Spółce na potrzeby nadzoru nad inwestycją, z zachowaniem poufności. Nie może przekazywać informacji niejawnych, Informacji Kontrolowanych ani informacji, których przekazanie wymaga zezwolenia na podstawie przepisów o kontroli eksportu. Do Dyrektora Inwestora stosuje się wyjątek portfelowy z \US{\ref{art:ip} ust.~5}.
\end{enumerate}
\end{ContractClause}

\begin{ContractClause}{Ust. 3 --- Wariant B: obserwator}
\label{upr:obserwator}
\begin{enumerate}
  \item Inwestor może wskazać jedną osobę uprawnioną do udziału w posiedzeniach Rady Dyrektorów bez prawa głosu (\textbf{„Obserwator”}) i w każdym czasie ją zmienić, w drodze oświadczenia w formie dokumentowej doręczonego Spółce.
  \item Obserwator otrzymuje zawiadomienia o posiedzeniach i materiały w tym samym czasie co dyrektorzy. Przed pierwszym posiedzeniem podpisuje umowę o zachowaniu poufności według wzoru Spółki, z zastosowaniem wariantu dla inwestora finansowego.
  \item Rada Dyrektorów może wyłączyć Obserwatora z części posiedzenia i nie przekazać mu materiałów, jeżeli dotyczą: konfliktu interesów Inwestora, relacji Spółki z Inwestorem, informacji objętych tajemnicą adwokacką albo radcowską, informacji niejawnych albo Informacji Kontrolowanych.
\end{enumerate}
\end{ContractClause}

\begin{ContractClause}{Ust. 4 --- Wariant C: prawa informacyjne}
\label{upr:informacje}
\begin{enumerate}
  \item Spółka przekazuje Inwestorowi w formie elektronicznej:
  \begin{enumerate}[label=\alph*)]
    \item raport miesięczny (przychody, koszty, środki pieniężne, zatrudnienie, kluczowe zdarzenia) --- w terminie 30 dni po upływie miesiąca;
    \item kwartalne sprawozdanie zarządcze --- w terminie 45 dni po upływie kwartału;
    \item roczne sprawozdanie finansowe, a jeżeli podlega badaniu --- ze sprawozdaniem z badania --- w terminie 14 dni od jego sporządzenia;
    \item budżet na następny rok obrotowy --- w terminie 14 dni od jego przyjęcia przez Radę Dyrektorów;
    \item kopię zgłoszeń do rejestru przedsiębiorców i wyciągu z rejestru akcjonariuszy --- na żądanie.
  \end{enumerate}
  \item Inwestor może raz w roku obrotowym, na własny koszt i po zawiadomieniu z wyprzedzeniem 14 dni, przejrzeć księgi i dokumenty Spółki przez siebie albo przez doradcę związanego tajemnicą zawodową.
  \item Prawa informacyjne nie obejmują informacji niejawnych, Informacji Kontrolowanych ani informacji, których przekazanie Inwestorowi wymagałoby zezwolenia na podstawie przepisów o kontroli eksportu albo naruszałoby warunki koncesji Spółki. Spółka nie przekazuje informacji podmiotowi, który jest konkurentem Spółki albo jest przez niego kontrolowany.
\end{enumerate}
\end{ContractClause}

\begin{ContractClause}{Ust. 5 --- Wariant D: zgody Inwestora}
\label{upr:zgody}
\begin{enumerate}
  \item Następujące czynności Spółki wymagają uprzedniej zgody Inwestora wyrażonej w formie dokumentowej, niezależnie od uchwał wymaganych umową spółki:
  \begin{enumerate}[label=\alph*)]
    \item zmiana umowy spółki pogarszająca prawa związane z akcjami serii objętej przez Inwestora albo uprzywilejowanie akcji innej serii w zakresie dywidendy, głosu albo podziału majątku;
    \item emisja akcji po cenie niższej niż cena akcji objętych przez Inwestora, z wyjątkiem akcji serii~F i~P;
    \item zwiększenie Puli programu motywacyjnego ponad \field{15}\% akcji Spółki po pełnym rozwodnieniu;
    \item rozwiązanie, likwidacja, połączenie, podział albo przekształcenie Spółki;
    \item zbycie przedsiębiorstwa albo jego zorganizowanej części albo przeniesienie kluczowej własności intelektualnej;
    \item zmiana przedmiotu działalności, która polega na rezygnacji z \field{obszar};
    \item \field{inne sprawy wynegocjowane w rundzie}.
  \end{enumerate}
  \item Zgodę uważa się za wyrażoną, jeżeli Inwestor nie odmówi jej w formie dokumentowej w terminie 10 dni roboczych od doręczenia wniosku wraz z niezbędnymi informacjami.
  \item Brak zgody nie wpływa na ważność czynności wobec osób trzecich; dyrektorzy odpowiadają wobec Spółki za dokonanie czynności bez zgody.
\end{enumerate}
\end{ContractClause}

\textbf{Zmiana \US{\ref{art:rada} ust.~2}} --- po słowach „uchwałą Walnego Zgromadzenia podjętą bezwzględną większością głosów oddanych” dodaje się: „, z zastrzeżeniem uprawnień indywidualnych, o których mowa w §~\field{numer}”.

% ================================================================
\chapter*{Część II --- Dokumenty wykonawcze}
\addcontentsline{toc}{chapter}{Część II --- Dokumenty wykonawcze}

\section*{Wzór uchwały Walnego Zgromadzenia}
\addcontentsline{toc}{section}{Wzór uchwały Walnego Zgromadzenia}

\textbf{Uchwała nr \field{numer} Walnego Zgromadzenia \CompanyFull{} z dnia \field{data} w sprawie zmiany umowy spółki --- uprawnienia indywidualne \field{firma inwestora}}

{\small (protokół notarialny; większość 75\% wszystkich głosów, \US{\ref{art:wz-zastrzezone} ust.~1})\par}

\begin{enumerate}
  \item Na podstawie art.~300$^{28}$ KSH Walne Zgromadzenie zmienia umowę spółki w ten sposób, że: (a)~po \US{\ref{art:wz-zastrzezone}} dodaje się paragraf „Uprawnienia indywidualne” w brzmieniu stanowiącym załącznik do uchwały; (b)~\US{\ref{art:rada} ust.~2} otrzymuje brzmienie: \field{brzmienie}.
  \item Walne Zgromadzenie upoważnia Radę Dyrektorów do ustalenia tekstu jednolitego umowy spółki.
  \item Uchwała wchodzi w życie z chwilą podjęcia, z mocą od dnia wpisu zmiany do rejestru. \field{Jeżeli zmiana jest elementem rundy:} Uchwała wchodzi w życie pod warunkiem wpisu do rejestru emisji akcji serii \field{oznaczenie} objętych przez Inwestora.
\end{enumerate}

\section*{Wzór oświadczenia o powołaniu albo odwołaniu Dyrektora Inwestora}
\addcontentsline{toc}{section}{Wzór oświadczenia o powołaniu albo odwołaniu Dyrektora Inwestora}

\field{firma Inwestora}, działając na podstawie art.~300$^{28}$ KSH i §~\field{numer} ust.~2 umowy spółki \CompanyFull{}, \field{powołuje / odwołuje} \field{imię, nazwisko, PESEL} \field{do / z} Rady Dyrektorów Spółki ze skutkiem od dnia \field{data}. \field{Przy powołaniu:} Do oświadczenia dołącza się zgodę powołanego na pełnienie funkcji, jego adres do doręczeń, oświadczenie o braku przeszkód z art.~18 KSH oraz informację, czy posiada poświadczenie bezpieczeństwa.

\field{miejscowość, data} \hfill \field{podpisy osób uprawnionych do reprezentacji Inwestora}

\section*{Lista kontrolna}
\addcontentsline{toc}{section}{Lista kontrolna}

\begin{enumerate}
  \item Warunki rundy mieszczą się w \UA{\ref{sha:runda} ust.~2}; jeżeli nie --- wymagana zmiana umowy akcjonariuszy.
  \item Ocena wymogu zawiadomienia na podstawie ustawy o kontroli niektórych inwestycji i warunków koncesji przed zawarciem umowy inwestycyjnej \field{do weryfikacji}.
  \item Uchwała WZ z protokołem notarialnym; zgłoszenie zmiany do rejestru z zaznaczeniem uprawnień indywidualnych (por. art.~300$^{12}$ §~2 pkt~5 KSH); tekst jednolity.
  \item Oświadczenie o powołaniu Dyrektora Inwestora albo wskazaniu Obserwatora; NDA Obserwatora; zgłoszenie dyrektora do rejestru.
  \item Uzupełnienie regulaminu Rady Dyrektorów o udział Obserwatora i zasady przekazywania informacji Inwestorowi.
\end{enumerate}

\end{document}
